% !TeX root = ../main.tex

% ==================================================
% 1.9. ML/DL 프로젝트 실행 구조 통합
% ==================================================

\subsection{ML/DL 프로젝트 실행 구조 통합}

\subsubsection{지금까지 배운 내용이 왜 필요했을까?}

먼저 우리가 최종적으로 실행하려는 명령을 살펴보자.

\begin{dhkzsh}
conda activate transformer

python -m src.train \
    --data ./data/train.csv \
    --epochs 10 \
    --batch-size 32 \
    --output ./checkpoints/model.pt
\end{dhkzsh}

이 명령에는 지금까지 학습한 거의 모든 개념이 포함되어 있다.

\begin{dhktable}
|명령 요소	|관련 학습 내용|
|:---|:---|
|conda activate transformer|	2장: Python 실행 환경|
|python	|2장: Python 인터프리터|
|-m src.train	|4장: 모듈과 패키지|
|--data	|3장: 커맨드라인 인수|
|./data/train.csv	|1장: 상대 경로와 CWD|
|--epochs 10	|3장: argparse와 타입 변환|
|--batch-size 32	|3장: 커맨드라인 옵션|
|--output	|3장: 실행 인수|
|./checkpoints/model.pt|	5장: 파일 입출력|
\end{dhktable}

그리고 명령에는 직접 드러나지 않지만, 내부적으로 다음 개념들도 사용하게 된다.

\begin{dhkcode}
import os
import sys
from pathlib import Path

sys.executable
sys.argv
sys.path

os.environ
os.getcwd()

Path.cwd()
Path(__file__)
\end{dhkcode}

여기에 이전에 배운 \zsh{subprocess}까지 연결하면,

여러 학습 실험을 자동으로 실행하는 구조도 만들 수 있다.

즉, 9장의 목표는 개별 Python 기능을 하나의 프로젝트 실행 시스템으로 구성하는 것이다.

\subsubsection{프로젝트의 기본 디렉터리 구조}

우선 다음과 같은 프로젝트를 가정하자.

\begin{dhkdir}[in=6]
transformer-project/
>src/
>>__init__.py
>>train.py
>>config.py
>>dataset.py
>>model.py
>data/
>>train.csv
>checkpoints/
>logs/
>configs/
>>train.json
>scripts/
>>run_experiments.py
>.env
>.gitignore
>environment.yml
>README.md
\end{dhkdir}

각 디렉터리와 파일의 역할을 살펴보자.

\begin{dhktable}
|구성 요소	|역할|
|:---|:---|
|src/	|Python 소스 코드|
|src/train.py|	모델 학습 진입점|
|src/config.py|	학습 설정 관리|
|src/dataset.py	|데이터 로딩과 전처리|
|src/model.py|	모델 정의|
|data/|	학습 데이터|
|checkpoints/	|모델 가중치 및 체크포인트|
|logs/|	학습 로그|
|configs/|	학습 설정 파일|
|scripts/|	실험 실행 및 자동화|
|.env|	로컬 환경변수 설정|
|.gitignore	|Git 추적 제외 규칙|
|environment.yml|	Conda 환경 정의|
|README.md|	프로젝트 실행 방법|
\end{dhktable}

여기서 중요한 것은 파일의 종류가 아니라 책임을 기준으로 디렉터리를 분리했다는 것이다.

예를 들어 모델의 구조를 변경하려면 \zsh{model.py},

데이터 전처리를 변경하려면 \zsh{dataset.py},

학습 실행 옵션을 변경하려면 \zsh{train.py} 또는 설정 파일을 수정하는 식이다.

물론 실제 프로젝트의 규모에 따라 이 구조는 단순해질 수도 있고 복잡해질 수도 있다.

이번 학습에서는 이 구조를 기준으로 진행한다.

\subsubsection{프로젝트 루트(Project Root)와 CWD}

이번 장에서 가장 먼저 확실히 구분해야 할 개념이 있다.

바로 \zsh{프로젝트 루트와 현재 작업 디렉터리}이다.

우리는 \zsh{Path.cwd()}와 \zsh{Path(__file__)}의 차이를 배웠다.

그런데 실제 프로젝트에서는 이 차이가 매우 중요해진다.

\begin{dhknote}
\begin{dhksmry}[centered]
프로젝트 루트란?
\end{dhksmry}
\end{dhknote}

프로젝트 루트는 프로젝트의 최상위 디렉터리를 의미한다.

앞의 예제에서는 다음 디렉터리가 프로젝트 루트다.

\begin{dhkzsh}
transformer-project/
\end{dhkzsh}

여기에 \zsh{src, data, checkpoints} 등이 포함되어 있다.

하지만 주의해야 할 점이 있다.

\textbf{프로젝트 루트는 Python이나 운영체제가 자동으로 지정하는 특별한 디렉터리가 아니다}.

개발자가 프로젝트의 구조를 설계하면서 정하는 논리적인 기준 디렉터리다.

반면 CWD는 운영체제가 각 프로세스에 대해 관리하는 실행 환경의 일부이다.

따라서 다음 두 개념은 구분해야 한다.

\begin{dhktable}
|개념	|의미|
|:---|:---|
|Project Root|	개발자가 정의한 프로젝트 기준 디렉터리|
|CWD|	현재 프로세스의 작업 디렉터리|
\end{dhktable}

두 경로가 같을 수도 있지만 반드시 같아야 하는 것은 아니다.

\begin{dhknote}
\begin{dhksmry}[centered]
같은 코드를 다른 위치에서 실행하면?
\end{dhksmry}
\end{dhknote}

프로젝트 구조가 다음과 같다고 가정하자.

\begin{dhkdir}[in=8]
/Users/user/Desktop/transformer-project/
>src/
>>__init__.py
>>train.py
>data/
>train.csv
\end{dhkdir}

\zsh{train.py}에 다음 코드를 작성해 보자.

\begin{dhkcode}
from pathlib import Path

print("CWD:", Path.cwd())
print("File:", Path(__file__).resolve())
\end{dhkcode}

먼저 프로젝트 루트에서 실행해 보자.

\begin{dhkzsh}
cd /Users/user/Desktop/transformer-project

python -m src.train
\end{dhkzsh}

결과:

\begin{dhkzsh}
CWD: /Users/user/Desktop/transformer-project
File: /Users/user/Desktop/transformer-project/src/train.py
\end{dhkzsh}

이번에는 상위 디렉터리에서 실행해 보자.

\begin{dhkzsh}
cd /Users/user/Desktop

python transformer-project/src/train.py
\end{dhkzsh}

결과:

\begin{dhkzsh}
CWD: /Users/user/Desktop
File: /Users/user/Desktop/transformer-project/src/train.py
\end{dhkzsh}

여기서 무엇이 달라졌을까?

\zsh{Path.cwd()}는 변경되었지만, \zsh{Path(__file__).resolve()}는 동일한 파일을 가리키고 있다.

즉,

\begin{dhkcode}
Path.cwd()
\end{dhkcode}

는 현재 프로세스가 어디에서 실행되고 있는지를 나타내고,

\begin{dhkcode}
Path(__file__).reslove()
\end{dhkcode}

는 현재 소스 파일의 경로를 나타내는 것이다.

단, \zsh{__file__}은 일반적인 Python 소스 모듈에서는 사용할 수 있지만

Jupyter Notebook이나 대화형 인터프리터에서는 정의되지 않을 수 있다.

\subsubsection{상대 경로는 무엇을 기준으로 해석될까?}

다음 코드를 보자.

\begin{dhkcode}
from pathlib import Path

data_path = Path("./data/train.csv")

print(data_path)
\end{dhkcode}

이 코드에서 \zsh{data_path}는 상대 경로를 표현하는 \zsh{Path} 객체다.

그런데 실제 파일에 접근할 때는 이 경로를 어디를 기준으로 해석할까?

일반적인 파일 접근에서는 \zsh{현재 작업 디렉터리(CWD)}를 기준으로 해석한다.

예를 들어 프로젝트 루트에서 실행하면

\begin{dhkzsh}
CWD = /Users/user/Desktop/transformer-project
\end{dhkzsh}

다음 경로는

\begin{dhkcode}
Path("./data/train.csv")
\end{dhkcode}

실질적으로 다음 파일을 가리킨다.

\begin{dhkzsh}
/Users/user/Desktop/transformer-project/data/train.csv
\end{dhkzsh}

하지만 상위 디렉터리에서 실행하면

\begin{dhkzsh}
CWD = /Users/user/Desktop
\end{dhkzsh}

동일한 상대 경로가 다음 파일을 가리키게 된다.

\begin{dhkzsh}
/Users/user/Desktop/data/train.csv
\end{dhkzsh}

따라서 파일을 찾지 못하는 문제가 발생할 수 있다.

이것이 프로젝트 실행 위치에 따라 \zsh{FileNotFoundError}가 발생하는 대표적인 이유 중 하나다.

\begin{dhknote}
\begin{dhksmry}[centered]
해결 방법 : 프로젝트 루트 기준 경로
\end{dhksmry}
\end{dhknote}

프로젝트 루트를 코드에서 명시적으로 계산할 수도 있다.

\begin{dhkcode}
from pathlib import Path

PROJECT_ROOT = Path(__file__).resolve().parent.parent

DATA_PATH = PROJECT_ROOT / "data" / "train.csv"

print(DATA_PATH)
\end{dhkcode}

여기서 경로를 단계별로 살펴보자.

\begin{dhkcode}
Path(__file__).resolve()
\end{dhkcode}

결과:

\begin{dhkzsh}
/Users/user/Desktop/transformer-project/src/train.py
\end{dhkzsh}

다음으로

\begin{dhkcode}
Path(__file__).resolve().parent
\end{dhkcode}

결과:

\begin{dhkzsh}
/Users/user/Desktop/transformer-project/src
\end{dhkzsh}

그리고

\begin{dhkcode}
Path(__file__).resolve().parent.parent
\end{dhkcode}

결과:

\begin{dhkzsh}
/Users/user/Desktop/transformer-project
\end{dhkzsh}

따라서 \zsh{PROJECT_ROOT}를 얻을 수 있다.

이제 다음 코드로 데이터를 참조하면

\begin{dhkcode}
DATA_PATH = PROJECT_ROOT / "data" / "train.csv"
\end{dhkcode}

현재 작업 디렉터리와 무관하게 동일한 프로젝트 파일을 가리킬 수 있다.

다만 이 방식은 \zsh{train.py}가 프로젝트 루트 바로 아래의 \zsh{src/}에 위치한다는 구조적 가정에 의존한다.

프로젝트 구조가 바뀌면 \zsh{parent.parent}도 적절히 변경해야 한다.

\subsubsection{CLI로 전달한 상대 경로는 어떻게 처리할까?}

여기서 이전에 배운 \zsh{argparse}와 연결해 보자.

다음 코드를 작성했다고 가정하자.

\begin{dhkcode}
import argparse
from pathlib import Path

parser = argparse.ArgumentParser()

parser.add_argument(
    "--data",
    type=Path,
    required=True
)

args = parser.parse_args()

print(args.data)
\end{dhkcode}

터미널에서 다음과 같이 실행하면

\begin{dhkzsh}
python -m src.train --data ./data/train.csv
\end{dhkzsh}

\zsh{args.data}에는 어떤 객체가 저장될까?

\begin{dhkcode}
Path("data/train.csv")
\end{dhkcode}

와 같은 상대 경로를 표현하는 \zsh{Path} 객체가 저장된다.

여기서 중요한 점이 있다.

\zsh{argparse}의 \zsh{type=Path}는 문자열을 \zsh{Path} 객체로 반환할 뿐,

자동으로 프로젝트 루트 기준의 절대 경로를 만들어 주지는 않는다.

따라서 다음 코드는

\begin{dhkcode}
args.data.is_file()
\end{dhkcode}

현재 작업 디렉터리를 기준으로 파일 존재 여부를 확인하게 된다.

\begin{dhknote}
\begin{dhksmry}[centered]
CLI 경로를 처리하는 두 가지 설계
\end{dhksmry}
\end{dhknote}

실제 프로젝트에서는 크게 두 가지 방식이 있다.

방식 A - 사용자의 CWD를 기준으로 해석

\begin{dhkcode}
data_path = args.data.resolve()
\end{dhkcode}

이 방식은 사용자가 터미널에서 전달한 상대 경로를 현재 작업 디렉터리 기준으로 해석한다.

예를 들어 사용자가

\begin{dhkzsh}
python -m src.train --data ./data/train.csv
\end{dhkzsh}

를 실행했다면, 사용자가 명령을 실행한 위치를 기준으로 파일을 찾는 것이다.

일반적인 CLI 프로그램에서는 자연스로운 방식이다.

방식 B - 프로젝트 루트를 기준으로 해석

\begin{dhkcode}
data_path = args.data

if not data_path.is_absolute():
    data_path = PROJECT_ROOT / data_path
\end{dhkcode}

이 방식은 상대 경로를 프로젝트 루트 기준으로 해석한다.

프로젝트 내부 데이터와 설정 파일을 일관된 위치에서 관리하려는 경우 유용하다.

하지만 사용자가 다른 디렉터리에 있는 파일을 상대 경로로 전달했을 때 예상과 다르게 동작할 수도 있다.

따라서 어느 방식이 절대적으로 옳다고 할 수는 없다.

중요한 것은 CLI가 받은 상대 경로를 어떤 기준으로 해석할지 명시적으로 설계하는 것이다.

이번 프로젝트에서는 학습 명령을 프로젝트 루트에서 실행하는 것을 기본 규칙으로 삼고, CLI 경로는 CWD 기준으로 처리한다.

어이렇게 하면 일반적인 Python CLI의 동작 방식과도 일치한다.

\subsubsection{\zsh{python train.py}와 \zsh{python -m src.train}}

이번에는 이전에 배웠던 import와 패키지 구조를 연결해 보자.

다음 프로젝트가 있다고 가정하자.

\begin{dhkdir}[in=6]
transformer-project/
>src/
>>__init__.py
>>train.py
>>dataset.py
>>model.py
\end{dhkdir}

\zsh{train.py}에서 같은 패키지에 있는 \zsh{dataset.py}를 가져오고 싶다.

\begin{dhkcode}
from .dataset import load_dataset
\end{dhkcode}

여기서 \zsh{.}은 현재 패키지를 기준으로 하는 상대 import를 의미한다.

그런데 다음처럼 실행하면 문제가 발생할 수 있다.

\begin{dhkzsh}
python src/train.py
\end{dhkzsh}

왜 그럴까?

Python은 \zsh{train.py}를 패키지 내부의 모듈이 아니라 직접 실행한 스크립트로 취급하기 때문이다.

이 경우 일반적으로

\begin{dhkzsh}
__name__="__main__"
\end{dhkzsh}

이고,

\begin{dhkzsh}
__package__ is None
\end{dhkzsh}

이므로 상대 import를 해석할 패키지 정보가 없다.

반면 다음과 같이 실행하면

\begin{dhkzsh}
python -m src.train
\end{dhkzsh}

Python은 \zsh{src.train}을 패키지 구조 안에서 찾아 실행한다.

이때 실행되는 모듈의 정보는 일반적으로 다음과 같다.

\begin{dhkzsh}
__name__=="__main__"

__package__=="src"
\end{dhkzsh}

따라서 다음 상대 import가 가능해진다.

\begin{dhkcode}
from .dataset import load_dataset
\end{dhkcode}

이를 정리하면 다음과 같다.

\begin{dhktable}
|실행 방식	|__name__|	__package__	|상대 import|
|:---|:---|:---|:---|
|python src/train.py|	"__main__"	|일반적으로 None	|일반적으로 실패|
|python -m src.train|"__main__"|	"src"|	가능|
\end{dhktable}

물론 \zsh{python -m src.train}이 성공하려면 

\zsh{src} 패키지를 Python이 모듈 검색 경로에서 찾을 수 있어야 한다.

이번 예제에서는 프로젝트 루트에서 실행하기 때문에 해당 조건을 만족한다.

왜 프로젝트에서는 \zsh{-m} 실행을 권장할까?

핵심은 소스 파일의 물리적 경로가 아니라 패키지 구조를 기준으로 실행할 수 있기 때문이다.

프로젝트가 커질수록 모듈 사이의 import 관계가 복잡해지는데, 이때 패키지 구조를 일관되게 유지하는 것이 중요하다.

\subsubsection{학습 진입점(Entry Point) 설계}

이제 실제 \zsh{train.py}를 구성해 보자.

먼저 학습 진입점이라는 개념을 이해해야 한다.

Entry Point는 프로그램 실행이 시작되는 진입 위치를 의미한다.

이번 프로젝트에서는 다음 명령이 학습 진입점을 실행한다.

\begin{dhkzsh}
python -m src.train
\end{dhkzsh}

\zsh{src/train.py}를 다음처럼 작성해 보자.

\begin{dhkcode}
# src/train.py

import argparse
import os
import sys
from pathlib import Path


def parse_args():
    parser = argparse.ArgumentParser(
        description="ML training script"
    )

    parser.add_argument(
        "--data",
        type=Path,
        required=True
    )

    parser.add_argument(
        "--epochs",
        type=int,
        default=10
    )

    parser.add_argument(
        "--batch-size",
        type=int,
        default=32
    )

    parser.add_argument(
        "--output",
        type=Path,
        default=Path("checkpoints/model.pt")
    )

    return parser.parse_args()


def main():
    args = parse_args()

    print("Python:", sys.executable)
    print("CWD:", Path.cwd())
    print("PID:", os.getpid())

    print("Data:", args.data)
    print("Epochs:", args.epochs)
    print("Batch size:", args.batch_size)
    print("Output:", args.output)


if __name__ == "__main__":
    main()
\end{dhkcode}

이 구조를 자세히 살펴보자.

\begin{dhknote}
\begin{dhksmry}[centered]
parse_args()를 분리하는 이유
\end{dhksmry}
\end{dhknote}

다음 코드를 별도의 함수로 분리했다.

\begin{dhkcode}
def parse_args():
    ...
\end{dhkcode}

이유는 커맨드라인 인수 처리와 학습 실행 로직의 책임을 구분하기 위해서다.

예를 들어 추후에 학습 설정을 검증하거나 테스트 코드를 작성할 때 구조를 확장하기 쉬워진다.

또한 더 큰 프로젝트에서는 인수 파싱 자체를 별도의 모듈로 분리할 수도 있다.

\begin{dhknote}
\begin{dhksmry}[centered]
main()을 사용하는 이유
\end{dhksmry}
\end{dhknote}

다음 함수는 프로그램의 주요 실행 흐름을 담당한다.

\begin{dhkcode}
def main():
    ...
\end{dhkcode}

Python에서 \zsh{main()}이라는 이름은 특별한 예약어가 아니다.

개발자가 관례적으로 사용하는 함수 이름이다.

즉, Python이 자동으로 \zsh{main()}을 실행해 주는 것은 아니다.

따라서 다음 코드가 필요하다.

\begin{dhkcode}
if __name__ == "__main__":
    main()
\end{dhkcode}

이 코드는 현재 모듈이 이 프로그램의 진입점으로 실행되었을 때 \zsh{main()}을 호출한다.

반면 다른 모듈에서 다음처럼 import한다면

\begin{dhkcode}
from src.train import main
\end{dhkcode}

\zsh{main()}이 자동으로 실행되지는 않는다.

이것은 프로젝트를 모듈 단위로 재사용하거나 테스트할 때 중요한 설계다.

\subsubsection{입력 파일과 출력 파일의 경로 검증}

이제 실제 학습 프로그램에서 필요한 경로 처리를 추가해 보자.

먼저 입력 파일을 확인해야 한다.

\begin{dhkcode}
if not args.data.is_file():
    raise FileNotFoundError(
        f"Dataset not found: {args.data}"
    )
\end{dhkcode}

입력 파일이 존재하지 않으면 학습을 시작하기 전에 오류를 발생시키는 것이다.

이번에는 출력 경로를 살펴보자.

\begin{dhkzsh}
args.output
\end{dhkzsh}

이 객체는 다음과 같은 경로를 가리킬 수 있다.

\begin{dhkzsh}
checkpoints/model.pt
\end{dhkzsh}

따라서 다음처럼 작성하여야 한다.

\begin{dhkcode}
args.output.parent.mkdir(
    parents=True,
    exist_ok=True
)
\end{dhkcode}

이전에 배웠던 내용이다.

\begin{dhkcode}
args.output.mkdir()
\end{dhkcode}

라고 작성하면 \zsh{model.pt}라는 이름의 디렉터리를 만들려고 하기 때문에 의도와 다른 동작이 발생한다.

\begin{dhknote}
\begin{dhksmry}[centered]
입력 경로와 출력 경로 처리의 차이
\end{dhksmry}
\end{dhknote}

\begin{dhktable}
|경로	|처리|
|:---|:---|
|입력 데이터 파일	|존재 여부 검증|
|출력 모델 파일	|부모 디렉터리 생성|
|로그 디렉터리	|디렉터리 자체 생성|
\end{dhktable}

예를 들어 다음처럼 구분할 수 있다.

\begin{dhkcode}
data_path = args.data.resolve()

output_path = args.output.resolve()

log_dir = Path("logs").resolve()

if not data_path.is_file():
    raise FileNotFoundError(data_path)

output_path.parent.mkdir(
    parents=True,
    exist_ok=True
)

log_dir.mkdir(
    parents=True,
    exist_ok=True
)
\end{dhkcode}

여기서 \zsh{resolve()}는 경로를 절대 경로로 해석하는 데 사용된다.

단, \zsh{resolve()} 자체가 파일을 생성하거나 파일 존재를 검증하는 것은 아니다.

파일 존재 여부는 별도로 확인해야 한다.

\subsubsection{실행 설정과 프로그램 로직의 분리}

이제 ML/DL 프로젝트에서 중요한 설계 원칙을 하나 살펴보자.

바로 설정(Configuration)과 실행 로직(Logic)의 분리다.

예를 들어 다음과 같이 학습 코드를 작성했다고 생각해 보자.

\begin{dhkzsh}
epochs = 10
batch_size = 32
learning_rate = 0.0001
\end{dhkzsh}

이 값들을 소스 코드 안에 직접 작성할 수도 있다.

하지만 실험할 때마다 코드를 수정해야 한다는 문제가 발생한다.

그래서 일반적으로 다음 세 가지 방법을 사용한다.

\begin{dhksmry}
    1. 커맨드라인 인수

    2. 환경변수

    3. 설정 파일
\end{dhksmry}

각각의 역할을 비교해 보자.

\begin{dhktable}
|방식|	적합한 용도	|예시|
|:---|:---|:---|
|CLI|	실행마다 변경하는 //*옵션	|epochs, batch size|
|환경변수|	실행 환경에 따른 설정	|API token, cache path|
|설정 파일|	여러 설정을 묶어서 //*관리	|optimizer, scheduler, model architecture|
\end{dhktable}

물론 엄격하게 구분되는 것은 아니다.

예를 들어 learning rate를 CLI로 전달할 수도 있고, JSON 설정 파일에 저장할 수도 있다.

중요한 것은 설정값의 책임과 변경 주기를 고려해 적절한 위치에 배치하는 것이다.

\begin{dhknote}
\begin{dhksmry}[centered]
JSON 설정 파일
\end{dhksmry}
\end{dhknote}

다음 파일을 만들어 보자.

\begin{dhkzsh}
configs/train.json
\end{dhkzsh}

내용은 다음과 같다.

\begin{dhkzsh}
{
    "model_name": "distilbert-base-uncased",
    "learning_rate": 0.0001,
    "weight_decay": 0.01,
    "max_length": 128
}
\end{dhkzsh}

Python에서는 다음과 같이 읽을 수 있다.

\begin{dhkcode}
import json
from pathlib import Path

config_path = Path("configs/train.json")

with config_path.open(
    "r",
    encoding="utf-8"
) as f:
    config = json.load(f)

print(config)
\end{dhkcode}

여기서 \zsh{json.load(f)}는 JSON 파일의 내용을 Python 객체로 변환한다.

결과는 다음과 같은 딕셔너리다.

\begin{dhkzsh}
{
    "model_name": "distilbert-base-uncased",
    "learning_rate": 0.0001,
    "weight_decay": 0.01,
    "max_length": 128
}
\end{dhkzsh}

이전에 배운 매핑 객체의 개념이 여기서 다시 등장한다.

\begin{dhkzsh}
config['learning_rate']
\end{dhkzsh}

는 \zsh{'learning_rate'} 키에 대응하는 값을 조회하는 연산이다.

\begin{dhknote}
\begin{dhksmry}[centered]
CLI와 설정 파일을 결합하기
\end{dhksmry}
\end{dhknote}

이번에는 설정 파일의 경로 자체를 CLI로 전달해 보자.

\begin{dhkcode}
parser.add_argument(
    "--config",
    type=Path,
    default=Path("configs/train.json")
)
\end{dhkcode}

이제 다음처럼 실행할 수 있다.

\begin{dhkzsh}
python -m src.train \
    --config configs/train.json \
    --epochs 10
\end{dhkzsh}

이 구조에서는 다음 두 종류의 설정을 구분할 수 있다.

\begin{dhkdir}[in=6]
CLI
> --config
>> 설정 파일 위치
>--epochs
>> 이번 실행의 epoch 수
\end{dhkdir}

설정 파일:

\begin{dhkdir}[in=6]
train.json
> model_name
> learning_rate
> weight_decay
> max_length
\end{dhkdir}

즉, CLI는 실행을 제어하고 설정 파일은 모델과 학습에 필요한 여러 설정을 묶어서 제공하는 구조다.

실제 프로젝트에서는 설정 파일과 CLI에 동일한 항목이 존재할 때 어느 쪽이 우선하는지도 명시적으로 설계해야 한다.

예를 들어 CLI 옵션이 설정 파일 값을 덮어쓰도록 만들 수 있다.

\subsubsection{지금까지의 실행 흐름 통합}

여기까지 배운 내용을 하나의 구조로 정리해 보자.

\begin{dhkbranch}[in=8]
node main :: Terminal
arrow main :: conda activate transformer
node main :: Conda Environment
arrow main :: Python Interpreter 선택
node main :: Shell
line main :: python -m src.train
line main ::   --data ./data/train.csv
line main ::   --epochs 10
arrow main ::   --output ./checkpoints/model.pt
node main :: Python Process
branchline main -> exe :: sys.executable
handoff exe -> exesub :: 현재 Python Interpreter
stop exesub
branchline main -> cwd :: Path.cwd()
handoff cwd -> cwdsub :: 현재 작업 디렉터리
stop cwdsub
branchline main -> argparse :: argparse
branchline argparse -> argdata :: data
stop argdata
branchline argparse -> argepoch :: epochs
stop argepoch
handoff argparse -> argoutput :: output
stop argoutput
branchline main -> import :: import
handoff import -> src :: src 패키지 모듈 로딩
stop src
branchline main -> json :: JSON Config
handoff json -> jsonsub :: 학습 설정 로딩
stop jsonsub
branchline main -> dataset :: Dataset
handoff dataset -> datasetsub :: data/train.csv
stop datasetsub
handoff main -> training :: Training
branchline training -> model :: Model
stop model
branchline training -> opt :: Optimizer
stop opt
branchline training -> check :: Checkpoint
stop check
handoff training -> logs :: Logs
stop logs
\end{dhkbranch}

이 구조에서 각 기능은 독립적인 책임을 가지고 있다.

그리고 중요한 것은 각 기능이 서로 다른 계층에 존재한다는 점이다.

예를 들어 다음 내용을 구분할 수 있어야 한다.

\begin{dhksmry}
- Shell은 명령을 해석하고 프로세스를 실행한다.

- Conda는 실행 환경을 구성하고 인터프리터 선택에 영향을 준다.

- Python 인터프리터는 모듈을 로딩하고 코드를 실행한다.

- argparse는 커맨드라인 인수를 Python 객체로 변환한다.

- pathlib는 파일 시스템 경로를 표현하고 조작한다.

- json은 설정 파일의 데이터를 Python 객체로 변환한다.

- 모델 학습 코드는 이 설정들을 이용해 실제 연산을 수행한다.
\end{dhksmry}

이 구분이 명확해지면 프로젝트에서 문제가 발생했을 때 어느 계층을 확인해야 하는지도 판단하기 쉬워진다.

\begin{dhknote}
\begin{dhksmry}[centered]
이번 학습 내용 정리
\end{dhksmry}
\end{dhknote}

\begin{dhktable}
|학습 항목|	핵심|
|:---|:---|
|프로젝트 루트|	프로젝트 구조의 논리적 기준|
|CWD	|현재 프로세스의 작업 디렉터리|
|Path(__file__)	|소스 파일의 경로|
|상대 경로	|일반적으로 CWD 기준으로 해석|
|python -m	|패키지 구조를 유지하며 모듈 실행|
|argparse	|실행 옵션을 Python 객체로 변환|
|main()	|프로그램의 실행 흐름을 구성하는 관례적 함수|
|입력 경로 검증|	is_file()|
|출력 경로 준비	|output.parent.mkdir()|
|JSON 설정|	json.load()|
|설정 분리	|CLI / 환경변수 / 설정 파일|
\end{dhktable}

\subsubsection{전체 학습 파이프 라인}

지난 학습에서는 프로젝트 루트, CWD, \zsh{python -m}, \zsh{argparse},

설정파일, 입력·출력 경로를 하나의 프로젝트 구조로 연결했다.

이번에는 그 구조 안에서 데이터가 어떻게 모델에 전달되고,

학습 결과가 어떻게 파일로 저장되는지 살펴볼 것이다.

여기서 중요한 것은 모델의 성능이 아니다.

이번에는 Python 프로그램의 실행 구조와 객체·데이터의 이동에 집중한다.

먼저 우리가 구현할 전체 흐름을 살펴본다.

\begin{dhkbranch}[in=8]
node main :: Terminal
line main ::python -m src.train
line main ::--data ./data/train.csv
line main :: --epochs 10
line main :: --batch-size 32
arrow main ::--output ./checkpoints/model.pt
node main :: src/train.py

branchline main -> argp :: argparse
handoff argp -> argpsub :: CLI 인수 파싱
stop argpsub

branchline main -> dataset :: dataset.py
handoff dataset -> datasetsub :: CSV → DataFrame → Tensor
stop datasetsub

branchline main -> model :: model.py
handoff model -> modelsub :: PyTorch Model 생성
stop modelsub

branchline main -> device :: Device
handoff device -> devicesub :: CPU / CUDA / MPS
stop devicesub

branchline main -> training :: Training Loop
branchline training -> ff :: Forward
stop ff
branchline training -> loss :: Loss
stop loss
branchline training -> bb :: Backward
stop bb
handoff training -> opt :: Optimizer Step
stop opt

handoff main -> check :: Checkpoint
handoff check -> checksub :: model.pt
stop checksub
\end{dhkbranch}

여기서 눈여겨볼 점은 각 단계가 서로 다른 역할을 수행한다는 것이다.

\begin{dhksmry}
    - argparse는 실행 설정을 전달받는다.

    - dataset.py는 외부 데이터를 모델이 사용할 수 있는 형태로 변환한다.

    - model.py는 학습할 모델을 정의한다.

    - train.py는 이 구성 요소들을 조립하고 실행한다.

    - torch.save()는 학습 결과를 파일로 저장한다.
\end{dhksmry}

이번에는 이해하기 쉽도록 CSV 기반의 간단한 회귀 모델을 사용한다.

Transformer 모델을 사용하면 토크나이저, 모델 다운롣, Hugging Face 캐시 등

추가적인 개념이 동시에 등장하기 때문이다.

그 부분은 기본 실행 구조를 완성한 뒤 연결할 것이다.

\subsubsection{학습 프로젝트 구성}

지난 학습에서 만든 프로젝트 구조를 그대로 사용한다.

\begin{dhkdir}[in=6]
transformer-project/
>src/
>>__init__.py
>>train.py
>>config.py
>>dataset.py
>>model.py
>data/
>>train.csv
>checkpoints/
>logs/
>configs/
>>train.json
>scripts/
>>run_experiments.py
>.env
>.gitignore
>environment.yml
>README.md
\end{dhkdir}

이번 예제에서 실제로 사용할 파일은 네 개다.

\begin{dhkzsh}
data/train.csv
src/dataset.py
src/model.py
src/train.py
\end{dhkzsh}

우선 데이터부터 살펴본다.

\begin{dhknote}
\begin{dhksmry}[centered]
CSV 데이터
\end{dhksmry}
\end{dhknote}

다음과 같은 데이터가 있다고 가정하자.

\begin{dhkzsh}
data/train.csv
\end{dhkzsh}

\begin{dhkzsh}
x1,x2,target
1.0,2.0,5.0
2.0,3.0,8.0
3.0,4.0,11.0
4.0,5.0,14.0
5.0,6.0,17.0
6.0,7.0,20.0
7.0,8.0,23.0
8.0,9.0,26.0
\end{dhkzsh}

이 데이터는 두 개의 입력 피처와 하나의 타깃을 가지고 있다.

수학적으로 표현하면 다음과 같다.

\[
X \in \mathbb{R}^{N\times 2}
\]
\[
y \in \mathbb{R}^{N\times 1}
\]

여기서 \(N\)은 데이터 샘플 수다.

우리는 간단한 선형 회귀 모델을 학습할 것이다.

\[
\hat{y}=XW+b
\]

이때

\[
W\in\mathbb{R}^{2\times1},
\qquad
b\in\mathbb{R}
\]

이다.

이 예제는 머신러닝 이론을 새롭게 학습하기 위한 것이 아니라,

데이터 파일이 Python 객체와 Tensor로 변환되는 과정을 확인하기 위한 것이다.

\subsubsection{\zsh{dataset.py} - 파일에서 Tensor까지}

이제 데이터 로딩 코드를 작성해 보자.

\begin{dhkcode}
# src/dataset.py

import pandas as pd
import torch

from pathlib import Path
from torch.utils.data import TensorDataset

def load_dataset(data_path: Path) -> TensorDataset:

    df = pd.read_csv(data_path)

    features = df[["x1", "x2"]]
    target = df[["target"]]

    x = torch.tensor(
        features.to_numpy(),
        dtype=torch.float32
    )

    y = torch.tensor(
        target.to_numpy(),
        dtype=torch.float32
    )

    dataset = TensorDataset(x, y)

    return dataset
\end{dhkcode}

이 코드에는 우리가 지금까지 배운 개념들이 많이 포함되어 있다.

하나씩 살펴보자.

\begin{dhknote}
\begin{dhksmry}[centered]
data_path: Path의 의미
\end{dhksmry}
\end{dhknote}

함수 정의를 보자.

\begin{dhkcode}
def load_dataset(data_path: Path) -> TensorDataset:
\end{dhkcode}

여기서 \zsh{data_path: Path}는 타입 힌트다.

이 함수가 \zsh{Path} 객체를 인수로 받도록 설계되었다는 것을 표현한다.

하지만 타입 힌트 자체가 런타임에서 인수의 타입을 강제로 검사하거나 변환하지는 않는다.

따라서 다음과 같이 호출할 수 있다.

\begin{dhkcode}
from pathlib import Path

dataset = load_dataset(
    Path("data/train.csv")
)
\end{dhkcode}

이때 \zsh{Path("data/train.csv")}는 파일 자체가 아니라 파일 시스템 경로를 표현하는 Python 객체다.

실제 파일을 읽는 작업은 다음 코드에서 발생한다.

\begin{dhkcode}
df = pd.read_csv(data_path)
\end{dhkcode}

이전에 배운 \zsh{Path} 객체와 파일 입출력의 차이가 여기서 다시 등장한다.

\begin{dhknote}
\begin{dhksmry}[centered]
pd.read_csv()는 무엇을 반환할까?
\end{dhksmry}
\end{dhknote}

\begin{dhkcode}
df = pd.read_csv(data_path)
\end{dhkcode}

\zsh{pd.read_csv()}는 CSV 데이터를 읽고 \zsh{DataFrame} 객체를 반환한다.

따라서 다음과 같은 바인딩이 발생한다.

\begin{dhkbranch}[in=8]
node main :: data_path
arrow main
node main :: Path Object
arrow main :: pd.read_csv()
node main :: DataFrame Object
arrowup main
node main :: df
\end{dhkbranch}

여기서 \zsh{df}는 파일 자체가 아니다.

파일의 내용을 읽어 구성한 메모리상의 \zsh{DataFrame} 객체를 참조하는 이름이다.

이 차이는 이후 체크포인트 저장에도 중요하게 등장한다.

\begin{dhknote}
\begin{dhksmry}[centered]
DataFrame에서 Tensor로 변환
\end{dhksmry}
\end{dhknote}

다음 코드를 살펴보자.

\begin{dhkcode}
features = df[["x1", "x2"]]
target = df[["target"]]
\end{dhkcode}

여기서 두 객체의 shape은 다음과 같다.

\begin{dhkzsh}
features.shape = (N, 2)
target.shape   = (N, 1)
\end{dhkzsh}

이제 NumPy 배열로 변환한다.

\begin{dhkcode}
features.to_numpy()
\end{dhkcode}

그리고 PyTorch Tensor로 변환한다.

\begin{dhkcode}
x = torch.tensor(
    features.to_numpy(),
    dtype=torch.float32
)
\end{dhkcode}

전체 흐름은 다음과 같다.

\begin{dhksmry}[centered]
CSV File
\da
pd.read_csv()
\da
DataFrame
\da
to_numpy()
\da
NumPy ndarray
\da
torch.tensor()
\da
PyTorch Tensor
\end{dhksmry}

여기서 \zsh{torch.tensor()}는 전달된 데이터를 이용해 새로운 Tensor를 생성한다.

또한 \zsh{dtype=torch.float32}를 지정했기 때문에

해당 Tensor의 자료형을 \zsh{float32}로 설정한다.

\subsubsection{\zsh{TensorDataset}과 \zsh{DataLoader}}

다음 코드를 다시 보자.

\begin{dhkcode}
dataset = TensorDataset(x, y)
\end{dhkcode}

\zsh{TensorDataset}은 여러 Tensor를 같은 인덱스 기준으로 묶어 주는 Dataset 구현체다.

예를 들어 다음처럼 접근할 수 있다.

\begin{dhkcode}
x_sample, y_sample = dataset[0]
\end{dhkcode}

이때 \zsh{x_sample}은 첫 번째 샘플의 피처를 나타내고,

\zsh{y_sample}은 첫 번째 샘플의 타깃을 나타낸다.

하지만 학습할 때는 일반적으로 데이터를 한 번에 모두 사용하기보다 미니배치 단위로 처리한다.

이때 \zsh{DataLoader}를 사용한다.

\begin{dhkcode}
from torch.utils.data import DataLoader

loader = DataLoader(
    dataset,
    batch_size=32,
    shuffle=True
)
\end{dhkcode}

이 코드를 객체 관점에서 살펴보자.

\begin{dhkbranch}[in=8]
node main :: TensorDataset Object
arrow main :: DataLoader(dataset, ...)
node main :: DataLoader Object
arrow main :: iter(loader)
node main :: Iterator
arrow main :: next(iterator)
node main :: Mini-batch
\end{dhkbranch}

여기서 이전에 배웠던 iterator와 Generator 개념이 연결된다.

정확하게 말하면 \zsh{DataLoader}는 iterable 객체이고,

\zsh{iter(loader)}를 호출하면 배치를 순회하는 iterator를 얻을 수 있다.

따라서 다음 코드가 가능하다.

\begin{dhkcode}
for x_batch, y_batch in loader:
    ...
\end{dhkcode}

이전에 배웠듯 \zsh{for}문은 내부적으로 iterator protocol을 사용한다.

개념적으로는 다음과 유사하다.

\begin{dhkcode}
iterator = iter(loader)

while True:
    try:
        x_batch, y_batch = next(iterator)

    except StopIteration:
        break
\end{dhkcode}

즉, \zsh{DataLoader}가 특별한 반복 문법을 사용하는 것이 아니다.

Python의 기존 iterator protocol을 활용하는 것이다.

이 부분은 이전에 학습한 Python 문법이 실제 ML/DL 코드에서 사용되는 대표적 사례다.

\subsubsection{\zsh{model.py} - 모델 객체 정의}

이제 모델을 정의해 보자.

\begin{dhkcode}
# src/model.py

import torch
from torch import nn

class RegressionModel(nn.Module):

    def __init__(self):
        super().__init__()

        self.linear = nn.Linear(2, 1)

    def forward(self, x):
        return self.linear(x)
\end{dhkcode}

여기서 모델은 다음과 같은 선형 변환을 수행한다.

\[
\hat{y}=XW+b
\]

다만 PyTorch의 \zsh{nn.Linear(2,1)}은 내부적으로 가중치를 (1,2) 형태로 저장하고

연산은 다음과 같이 수행한다.

\[
Y=XW^{T}+b
\]

이것은 행렬의 shape을 맞추기 위한 구현상의 표현 차이다.

\begin{dhknote}
\begin{dhksmry}[centered]
모델 객체 생성
\end{dhksmry}
\end{dhknote}

다음 코드를 보자.

\begin{dhkcode}
model = RegressionModel()
\end{dhkcode}

이 코드는 \zsh{RegressionModel} 클래스의 인스턴스를 생성한다.

\begin{dhkbranch}[in=8]
node main :: RegressionModel Class
arrow main :: RegressionModel()
node main :: RegressionModel Instance
branchline main -> lin :: linear
branchline lin -> weight :: weight
stop weight
handoff lin -> bias :: bias
stop bias
handoff main -> ff :: forward()
\end{dhkbranch}

여기서 \zsh{model}은 모델 파일이 아니다.

메모리에 생성된 Python 객체를 참조하는 이름이다.

이 모델 객체가 학습되면서 내부 파라미터의 값이 변경되는 것이다.

그리고 나중에 그 파라미터를 파일로 저장할 것이다.

\subsubsection{CPU, CUDA, MPS 장치 선택}

이제 모델을 어느 장치에서 실행할지 결정해야 한다.

사용하는 환경이 다양하다면 특히 중요한 부분이다.

\begin{dhksmry}
    - Apple Silicon Mac: MPS 사용 가능 여부 확인

    - Windows 데스크톱: CUDA 지원 GPU 여부 확인

    - 공통 대안: CPU
\end{dhksmry}

먼저 장치 선택 함수를 작성해 보자.

\begin{dhkcode}
import torch


def get_device():

    if torch.cuda.is_available():
        return torch.device("cuda")

    if torch.backends.mps.is_available():
        return torch.device("mps")

    return torch.device("cpu")
\end{dhkcode}

여기서 \zsh{torch.device()}는 연산을 수행할 장치를 표현하는 객체를 생성한다.

\begin{dhkcode}
device = get_device()
\end{dhkcode}

이후 모델을 해당 장치로 이동할 수 있다.

\begin{dhkcode}
model = model.to(device)
\end{dhkcode}

그리고 데이터도 같은 장치로 이동해야 한다.

\begin{dhkcode}
x_batch = x_batch.to(device)
y_batch = y_batch.to(device)
\end{dhkcode}

\begin{dhknote}
\begin{dhksmry}[centered]
왜 모델과 데이터를 같은 장치로 이동해야 할까?
\end{dhksmry}
\end{dhknote}

PyTorch의 일반적인 연산에서는 연산에 참여하는 Tensor들이 호환되는 장치에 있어야 한다.

예를 들어 모델의 가중치가 GPU에 있고, 입력 데이터가 CPU에 있으면 장치 불일치 오류가 발생할 수 있다.

\begin{dhkdir}[in=6]
CPU Memory
> Input Tensor
\end{dhkdir}

\begin{dhkdir}[in=6]
GPU Memory
> Model Parameters
\end{dhkdir}

이 상태에서는 일반적인 선형 계층 연산을 그대로 수행할 수 없다.

따라서 다음처럼 맞춰 준다.

\begin{dhkdir}[in=6]
GPU
> Model Parameters
> Input Tensor
> Target Tensor
\end{dhkdir}

\begin{dhknote}
\begin{dhksmry}[centered]
model.to(device)와 tensor.to(device)의 차이
\end{dhksmry}
\end{dhknote}

여기서 객체 관점에서 주의할 점이 있다.

\begin{dhkcode}
model.to(device)
\end{dhkcode}

\zsh{nn.Module.to()}는 일반적으로 모델 자체의 파라미터와 버퍼를 지정한 장치로 이동시키고, 모델 객체 자신을 반환한다.

반면

\begin{dhkcode}
x_batch.to(device)
\end{dhkcode}

Tensor의 \zsh{.to()}는 필요한 경우 변환된 Tensor를 반환한다.

따라서 Tensor에서는 반환값을 다시 바인딩하는 형태를 사용하는 것이 일반적이다.

\begin{dhkcode}
x_batch = x_batch.to(device)
\end{dhkcode}

이 차이는 Python 객체의 mutation과 rebinding을 공부했던 내용과도 연결된다.

\subsubsection{학습 루프 구성}

이제 실제 학습을 수행해 보자.

먼저 필요한 구성 요소를 생성한다.

\begin{dhkcode}
model = RegressionModel().to(device)

criterion = torch.nn.MSELoss()

optimizer = torch.optim.Adam(
    model.parameters(),
    lr=0.001
)
\end{dhkcode}

여기서 각각의 역할은 다음과 같다.

\begin{dhktable}
|객체	|역할|
|:---|:---|
|model	|예측값 계산|
|criterion	|손실 계산|
|optimizer|	모델 파라미터 갱신|
\end{dhktable}

학습 루프는 다음과 같이 구성할 수 있다.

\begin{dhkcode}
for epoch in range(epochs):

    model.train()

    for x_batch, y_batch in loader:

        x_batch = x_batch.to(device)
        y_batch = y_batch.to(device)

        optimizer.zero_grad()

        predictions = model(x_batch)

        loss = criterion(
            predictions,
            y_batch
        )

        loss.backward()

        optimizer.step()

    print(
        f"Epoch {epoch + 1}/{epochs}, "
        f"Loss: {loss.item():.4f}"
    )
\end{dhkcode}

이 학습 범위의 초점은 딥러닝의 수학적 원리가 아니므로 설명은 하지 않는다.

대신 Python 실행 구조 관점에서 살펴보자.

\begin{dhknote}
\begin{dhksmry}[centered]
model(x_batch)는 무엇일까?
\end{dhksmry}
\end{dhknote}

\begin{dhkcode}
predictions = model(x_batch)
\end{dhkcode}

이전에 callable 객체에 대해 배웠다.

\zsh{model}은 \zsh{RegressionModel}의 인스턴스이고,

\zsh{nn.Module}은 호출 가능한 객체로 동작한다.

따라서

\begin{dhkcode}
model(x_batch)
\end{dhkcode}

처럼 호출할 수 있다.

이 호출은 PyTorch의 \zsh{nn.Module} 호출 메커니즘을 거쳐 \zsh{forward()}를 실행한다.

여기서 중요한 점은 \zsh{model.forward(x_batch)}를 직접 호출하는 것보다

\zsh{model(x_batch)}를 사용하는 것이 일반적인 방식이라는 것이다.

PyTorch의 호출 메커니즘은 forward hook 등 추가적인 기능도 처리하기 때문이다.

\subsubsection{체크포인트(Checkpoint)란?}

이제 학습한 모델을 저장해 보자.

먼저 체크포인트라는 개념부터 구분해야 한다.

체크포인트는 특정 시점의 학습 상태를 저장한 데이터다.

예를 들어 다음 정보를 저장할 수 있다.

\begin{dhksmry}
    - 모델 파라미터

    - Optimizer 상태

    - 현재 Epoch

    - 학습 설정

    - 손실값
\end{dhksmry}

하지만 모든 체크포인트가 이 정보를 전부 포함하는 것은 아니다.

모델 파라미터만 저장하는 경우도 있고, 학습을 정확하게 이어가기 위한 추가 상태까지 저장하는 경우도 있다.

\begin{dhknote}
\begin{dhksmry}[centered]
모델 파라미터만 저장하기
\end{dhksmry}
\end{dhknote}

가장 단순한 방법은 다음과 같다.

\begin{dhkcode}
torch.save(
    model.state_dict(),
    "checkpoints/model.pt"
)
\end{dhkcode}

여기서 중요한 메서드는 다음이다.

\begin{dhkcode}
model.state_dict()
\end{dhkcode}

이 메서드는 모델의 파라미터와 persistent buffer 등을 담은 딕셔너리 형태의 객체를 반환한다.

따라서 다음 두 개념을 구분해야 한다.

\begin{dhkbranch}[in=8]
node main :: Model Object
arrow main :: state_dict()
node main :: State Dictionary
arrow main :: torch.save()
node main :: model.pt
\end{dhkbranch}

여기서 \zsh{model.pt}는 모델 객체 자체가 아니라, 저장된 데이터를 담고 있는 파일이다.

즉, 메모리의 Python 객체와 디스크의 파일을 구분해야 한다는 5장의 내용이 다시 등장한다.

\begin{dhknote}
\begin{dhksmry}[centered]
저장 경로 준비하기
\end{dhksmry}
\end{dhknote}

다음 코드를 작성해 보자.

\begin{dhkcode}
from pathlib import Path

output_path = Path("checkpoints/model.pt")

output_path.parent.mkdir(
    parents=True,
    exist_ok=True
)

torch.save(
    model.state_dict(),
    output_path
)
\end{dhkcode}

여기서 \zsh{output_path.parent}는 \zsh{checkpoints} 디렉터리를 가리킨다.

따라서 모델 파일을 저장하기 전에 부모 디렉터리를 생성한다.

\subsubsection{체크포인트 불러오기}

이번에는 저장한 모델 파라미터를 불러와 보자.

\begin{dhkcode}
model = RegressionModel()

state_dict = torch.load(
    "checkpoints/model.pt",
    map_location="cpu",
    weights_only=True
)

model.load_state_dict(state_dict)
\end{dhkcode}

여기서 중요한 부분은 두 단계로 나뉜다는 것이다.

첫 번째:

\begin{dhkcode}
state_dict = torch.load(...)
\end{dhkcode}

파일에 저장된 데이터를 읽어 Python 객체로 복원한다.

두 번째:

\begin{dhkcode}
model.load_state_dict(state_dict)
\end{dhkcode}

복원한 파라미터 값을 모델 객체에 적용한다.

따라서 다음 구조가 성립한다.

\begin{dhkbranch}[in=8]
node main :: model.pt
arrow main :: torch.load()
node main :: State Dictionary
arrow main :: model.load_state_dict()
node main :: Model Object
\end{dhkbranch}

여기서 \zsh{torch.load()}만 실행한다고 기존 모델 객체가 자동으로 생성되고 구성되는 것은 아니다.

우리는 먼저 다음 코드를 실행했다.

\begin{dhkcode}
model = RegressionModel()
\end{dhkcode}

즉, 모델의 구조를 정의하는 Python 코드가 필요하다.

그 구조에 저장된 파라미터를 적용하는 방식이다.

\zsh{weights_only=True}는 체크포인트를 불러올 때 허용되는 역직렬화 대상을 제한하는 옵션이다.

그렇더라도 신뢰할 수 없는 체크포인트를 무조건 안전하다고 가정해서는 안 된다.

\subsubsection{학습 재개용 체크포인트}

이번에는 모델 파라미터뿐 아니라 Optimizer 상태도 저장해 보자.

\begin{dhkcode}
checkpoint = {
    "epoch": epoch,
    "model_state_dict": model.state_dict(),
    "optimizer_state_dict": optimizer.state_dict(),
    "loss": loss.item()
}

torch.save(
    checkpoint,
    "checkpoints/checkpoint.pt"
)
\end{dhkcode}

여기서 \zsh{checkpoint}는 Python 딕셔너리다.

각 키에 학습 상태를 저장하고 있다.

\begin{dhkdir}[in=6]
checkpoint
> epoch
> model_state_dict
> optimizer_state_dict
> loss
\end{dhkdir}

이제 다시 불러오면 다음과 같다.

\begin{dhkcode}
checkpoint = torch.load(
    "checkpoints/checkpoint.pt",
    map_location="cpu",
    weights_only=True
)

model.load_state_dict(
    checkpoint["model_state_dict"]
)

optimizer.load_state_dict(
    checkpoint["optimizer_state_dict"]
)

start_epoch = checkpoint["epoch"] + 1
\end{dhkcode}

이렇게 하면 저장된 Epoch 이후부터 학습을 이어갈 수 있는 기본 구조가 만들어진다.

다만 정확하게 동일한 학습 상태를 재현하려면 학습률 스케줄러, 난수 생성기 상태, 데이터 샘플링 상태 등 추가적인 정보가 필요할 수도 있다.

따라서 모델 파라미터만 저장하는 것과 학습을 완전히 재개하기 위한 체크포인트를 저장하는 것은 구분해야 한다.

\subsubsection{\zsh{train.py}에서 전체 구성 요소 조립하기}

이제 지금까지 작성한 코드를 하나로 연결해 보자.

다음 코드는 학습 실행 구조를 이해하기 위한 최소 예제이다.

\begin{dhkcode}
# src/train.py

import argparse
import sys

from pathlib import Path

import torch

from torch.utils.data import DataLoader

from .dataset import load_dataset
from .model import RegressionModel


def parse_args():

    parser = argparse.ArgumentParser()

    parser.add_argument(
        "--data",
        type=Path,
        required=True
    )

    parser.add_argument(
        "--epochs",
        type=int,
        default=10
    )

    parser.add_argument(
        "--batch-size",
        type=int,
        default=32
    )

    parser.add_argument(
        "--output",
        type=Path,
        default=Path("checkpoints/model.pt")
    )

    return parser.parse_args()


def get_device():

    if torch.cuda.is_available():
        return torch.device("cuda")

    if torch.backends.mps.is_available():
        return torch.device("mps")

    return torch.device("cpu")


def main():

    args = parse_args()

    data_path = args.data.resolve()
    output_path = args.output.resolve()

    if not data_path.is_file():
        raise FileNotFoundError(data_path)

    output_path.parent.mkdir(
        parents=True,
        exist_ok=True
    )

    device = get_device()

    dataset = load_dataset(data_path)

    loader = DataLoader(
        dataset,
        batch_size=args.batch_size,
        shuffle=True
    )

    model = RegressionModel().to(device)

    criterion = torch.nn.MSELoss()

    optimizer = torch.optim.Adam(
        model.parameters(),
        lr=0.001
    )

    print("Python:", sys.executable)
    print("Device:", device)
    print("Data:", data_path)

    for epoch in range(args.epochs):

        model.train()

        total_loss = 0.0

        for x_batch, y_batch in loader:

            x_batch = x_batch.to(device)
            y_batch = y_batch.to(device)

            optimizer.zero_grad()

            predictions = model(x_batch)

            loss = criterion(
                predictions,
                y_batch
            )

            loss.backward()

            optimizer.step()

            total_loss += loss.item() * len(x_batch)

        avg_loss = total_loss / len(dataset)

        print(
            f"Epoch {epoch + 1}/{args.epochs} "
            f"Loss: {avg_loss:.4f}"
        )

    torch.save(
        model.state_dict(),
        output_path
    )

    print("Model saved:", output_path)


if __name__ == "__main__":
    main()
\end{dhkcode}

이 코드는 다음과 같이 실행할 수 있다.

프로젝트 루트에서:

\begin{dhkzsh}
python -m src.train \
    --data ./data/train.csv \
    --epochs 10 \
    --batch-size 32 \
    --output ./checkpoints/model.pt
\end{dhkzsh}

실행 결과는 대략 다음과 같다.

\begin{dhkzsh}
Python: /path/to/conda/envs/transformer/bin/python
Device: cpu
Data: /path/to/project/data/train.csv

Epoch 1/10 Loss: ...
Epoch 2/10 Loss: ...
...
Epoch 10/10 Loss: ...

Model saved: /path/to/project/checkpoints/model.pt
\end{dhkzsh}

실제 손실값은 모델의 초기 파라미터와 학습 환경에 따라 달라진다.

\subsubsection{지금까지 배운 Python 개념이 어디에 사용되었을까?}

이제 전체 코드를 Python 기초 개념과 연결해 보자.

\begin{dhktable}
|Python 개념	|실제 사용 위치|
|:---|:---|
|Module / Package	|from .dataset import load_dataset|
|python -m	|src.train을 패키지 모듈로 실행|
|argparse|	학습 인수 처리|
|Path	|데이터·체크포인트 경로|
|Path.resolve()	|상대 경로를 절대 경로로 해석|
|Path.is_file()	|입력 데이터 검증|
|Path.mkdir()	|출력 디렉터리 생성|
|Function|	parse_args(), get_device(), main()|
|Class / Instance	|RegressionModel()|
|Callable	|model(x_batch)|
|Iterator Protocol	|for x_batch, y_batch in loader|
|Mapping	|model.state_dict()|
|File I/O	|pd.read_csv(), torch.save()|
|Object Binding	|model = RegressionModel()|
|Python Interpreter|	sys.executable|
\end{dhktable}

여기서 주목할 점은 ML/DL 라이브러리가 Python의 기본 객체 모델 위에서 동작한다는 사실이다.

예를 들어 \zsh{DataLoader}는 반복 가능한 객체이고, 모델은 호출 가능한 객체이며,

\zsh{state_dict()}는 매핑 형태의 데이터를 반환한다.

즉, 지금까지 배운 Python 문법과 객체 모델이 실제 딥러닝 프로젝트의 기반이 되는 것이다.

\subsubsection{모델 파일과 학습 환경은 별개다}

마지막으로 중요한 개념 하나를 짚고 넘어간다.

우리가 다음 코드를 실행했다고 생각해 보자.

\begin{dhkcode}
torch.save(
    model.state_dict(),
    "checkpoints/model.pt"
)
\end{dhkcode}

그러면 모델 파라미터가 파일에 저장된다.

하지만 이 파일에 다음 정보가 모두 포함되는 것은 아니다.

\begin{dhksmry}
    - Python 인터프리터의 위치

    - Conda 환경 전체

    - 설치된 라이브러리 버전

    - 학습 데이터

    - 모델 클래스의 Python 소스 코드

    - 프로젝트 디렉터리 구조
\end{dhksmry}

즉, 모델 파일 하나만 가지고 언제나 동일한 실행 환경을 복원할 수 있는 것은 아니다.

이를 구분하면 다음과 같다.

\begin{dhkdir}[in=6]
ML/DL Project
> Source Code
>> src/
>Dependencies
>>environment.yml
>Configuration
>>configs/train.json
>Dataset
>>data/train.csv
>Checkpoint
>>checkpoints/model.pt
\end{dhkdir}

각 구성 요소는 서로 다른 책임을 가지고 있다.

그래서 실제 ML/DL 프로젝트에서는 모델 파라미터뿐 아니라 콛, 환경, 설정, 데이터의 버전까지 함께 관리하는 것이 중요하다.

이 부분은 추후 실험 재현성과 모델 배포를 공부할 때 다시 등장하게 될 것이다.

\begin{dhknote}
\begin{dhksmry}[centered]
이번 학습 정리
\end{dhksmry}
\end{dhknote}

이번에는 다음 흐름을 실제 코드로 연결했다.

\begin{dhksmry}[centered]
CSV 
\da 
DataFrame 
\da 
Tensor
\da
DataLoader
\da
Model
\da
Training
\da
Checkpoint 
\end{dhksmry}

그리고 이 전체 과정은 다음 실행 명령에서 시작했다.

\begin{dhkzsh}
python -m src.train \
    --data ./data/train.csv \
    --epochs 10 \
    --batch-size 32 \
    --output ./checkpoints/model.pt
\end{dhkzsh}

이제 우리가 처음에 배웠던 경로, CLI, import, 객체, 반복자, 파일 입출력이

실제 학습 파이프라인 안에서 어떻게 연결되는지 확인하였다.

\subsubsection{Hugging Face의 \zsh{from_pretrained()} 이해하기}

여기서부터는 Hugging Face의 transformers, Conda 환경, 모델 캐시,

환경변수, 실험 자동화를 하나의 구조로 연결한다.

앞에서 작성한 PyTorch 회귀 모델은 로컬 CSV 파일을 읽고 학습하는 구조였다.

하지만 Hugging Face를 사용하면 모델과 토크나이저를 외부 저장소에서 가져오고,

다운로드한 파일을 캐시에 저장하며, 모델에 따라 인증 토큰도 필요해진다.

즉, 이번 학습에서는 다음 질문에 답하는 것이 목표이다.

\begin{dhksmry}[centered]
Python 프로그램이 외부 모델을 가져와 학습하고,

그 실행 환경과 결과를 재현 가능하게 관리하려면 어떤 구조가 필요할까?
\end{dhksmry}

먼저 아래의 코드를 살펴보자.

\begin{dhkcode}
from transformers import AutoTokenizer, AutoModel

model_name = "distilbert-base-uncased"

tokenizer = AutoTokenizer.from_pretrained(model_name)
model = AutoModel.from_pretrained(model_name)
\end{dhkcode}

우리는 이 코드를 실행하면 사전학습된 모델과 토크나이저를 사용할 수 있다는 것을 알 수 있다.

그런데 Python 실행 환경 관점에서는 조금 더 깊게 이해할 필요가 있다.

\begin{dhknote}
\begin{dhksmry}[centered]
from_pretrained()는 무엇을 하는가?
\end{dhksmry}
\end{dhknote}

다음 코드를 보자.

\begin{dhkcode}
model = AutoModel.from_pretrained(
    "distilbert-base-uncased"
)
\end{dhkcode}

이때 \zsh{from_pretrained()}는 단순히 모델 객체를 생성하는 함수가 아니다.

일반적으로 다음 작업을 수행한다.

\begin{dhksmry}
    1. 모델 식별자를 해석한다.

    2. 필요한 설정 파일과 가중치 파일을 찾는다.

    3. 로컬 캐시에 사용할 수 있는 파일이 있는지 확인한다.

    4. 필요한 파일이 없다면 Hugging Face Hub에서 다운로드한다.

    5. 설정을 이용해 모델 구조를 생성한다.

    6. 저장된 가중치를 모델에 적용한다.

    7. 초기화된 모델 객체를 반환한다.
\end{dhksmry}

개념적으로 표현하면 다음과 같다.

\begin{dhkbranch}[in=8]
node main :: AutoModel.from_pretrained(model_name)
arrow main
node main :: Model Identifier
arrow main
node main :: Hugging Face Hub
arrow main :: 필요한 파일 조회
node main :: Local Cache
branchline main -> conf :: config.json
stop conf
handoff main -> weight :: model weights
arrow weight
node weight :: Model 생성
arrow weight
node weight :: 가중치 적용
arrow weight
node weight :: Model Object
stop weight
\end{dhkbranch}

단, 이미 캐시에 필요한 파일이 있으면 다운로드 하지 않고 사용할 수 있다.

실제 파일 조회 및 다운로드 여부는 캐시 상태와 라이브러리 옵션에 따라 달라진다.

여기서 중요한 점은 다음이다.

\zsh{from_pretrained()}가 반환하는 것은

메모리에 생성된 Python 모델 객체이지, 다운로드한 파일 경로가 아니다.

이전에 배운 \zsh{pd.read_csv()}와 비교해 보자.

\begin{dhktable}
|코드	|외부 데이터	|반환 객체|
|:---|:---|:---|
|pd.read_csv(path)	|CSV 파일	|DataFrame|
|torch.load(path)	|체크포인트 파일	|역직렬화된 객체|
|AutoModel//*.from_pretrained(//*name)	|모델 설정과 가중치|	모델 객체|
|AutoTokenizer//*.from_pretrained(//*name)|토크나이저 설정과//*관련 파일|토크나이저 객체|
\end{dhktable}

외부 데이터를 읽어서 Python 객체를 구성한다는 점에서는 비슷하지만, 각 함수가 수행하는 구체적인 작업은 다르다.

\subsubsection{모델 식별자와 파일 경로의 차이}

다음 문자열을 보자.

\begin{dhkcode}
model_name = "distilbert-base-uncased"
\end{dhkcode}

이 문자열은 일반적으로 Hugging Face Hub의 모델 저장소를 식별하는 데 사용된다.

반면 다음 코드는 어떨까?

\begin{dhkcode}
model_path = "./models/distilbert"
\end{dhkcode}

이것은 로컬 파일 시스템 경로로 사용할 수 있는 문자열이다.

\zsh{from_pretrained()}는 모델 식별자뿐 아니라 적절한 구조를 가진 로컬 디렉터리도 입력받을 수 있다.

예를 들어 다음과 같은 디렉터리가 있다고 가정하자.

\begin{dhkdir}[in=6]
model/
>distilbert/
>>config.json
>>model.safetensors
>>tokenizer_config.json
>>tokenizer.json
>>vocab.txt
\end{dhkdir}

실제 파일 구성은 모델과 토크나이저 종류에 따라 달라질 수 있다.

이제 다음처럼 사용할 수 있다.

\begin{dhkcode}
from pathlib import Path
from transformers import AutoTokenizer, AutoModel

model_dir = Path("models/distilbert")

tokenizer = AutoTokenizer.from_pretrained(
    model_dir
)

model = AutoModel.from_pretrained(
    model_dir
)
\end{dhkcode}

여기서 구분해야 할 개념은 세 가지다.

\begin{dhktable}
|개념|	의미|
|:---|:---|
|모델 식별자	|Hub에서 모델 저장소를 식별하는 문자열|
|로컬 모델 디렉터리	|모델 파일이 저장된 파일 시스템 경로|
|모델 객체|	Python 프로세스의 메모리에 생성된 모델 인스턴스|
\end{dhktable}

이 구분은 이전에 공부했던 이름, 객체, 경로, 파일의 차이와 직접 연결된다.

\subsubsection{Hugging Face 캐시(Cache)의 역할}

이제 캐시를 살펴보자.

이전에 Hugging Face 모델을 사용하면서 파일 다운로드 오류나 토크나이저 파일 문제를 경험한 적이 있을 것이다.

그런 문제를 이해하려면 캐시의 구조를 알아야 한다.

\begin{dhknote}
\begin{dhksmry}[centered]
캐시가 필요한 이유
\end{dhksmry}
\end{dhknote}

모델을 실행할 때마다 수백 MB 또는 수 GB의 파일을 다운로드한다면 비효율적이다.

그래서 Hugging Face는 다운로드한 파일을 로컬 캐시에 저장해 재사용할 수 있다.

일반적인 흐름은 다움과 같다.

\begin{dhkbranch}[in=8]
node main :: 필요한 파일 확인
branchline main -> ca :: 캐시에 사용 가능한 파일 존재
arrow ca
node ca :: 로컬 파일 사용
handoff main -> nf :: 파일 없음
arrow nf
node nf :: Hugging Face Hub
arrow nf
node nf :: 다운로드
arrow nf
node nf :: Cache 저장
arrow nf
node nf :: 파일 사용
\end{dhkbranch}

역시ㅓ 캐시는 모델 객체 자체를 저장하는 공간이 아니다.

모델 객체를 구성하는 데 필요한 파일을 저장하는 공간이다.

따라서 Python 프로세스가 종료되더라도 캐시 파일은 일반적으로 디스크에 남아 있다.

반면 메모리에 생성된 모델 객체는 프로세스 종료와 함께 사라진다.

\begin{dhknote}
\begin{dhksmry}[centered]
기본 캐시 위치
\end{dhksmry}
\end{dhknote}

Hugging Face Hub의 일반적인 기본 캐시 위치는 다음과 같다.

macOS/Linux:

\begin{dhkzsh}
~/.cache/huggingface/hub/
\end{dhkzsh}

Windows:

\begin{dhkzsh}
C:\Users\<username>\.cache\huggingface\hub\
\end{dhkzsh}

다만 환경변수와 라이브러리 설정에 따라 실제 위치는 달라질 수 있다.

여기서 \zsh{~}는 사용자 홈 디렉터리를 나타내는 표현이다.

Python에서는 다음과 같이 확인할 수 있다.

\begin{dhkcode}
from pathlib import Path

print(Path.home())
\end{dhkcode}

그리고 일반적인 기본 경로를 구성하면 다음과 같다.

\begin{dhkcode}
cache_dir = (
    Path.home()
    / ".cache"
    / "huggingface"
    / "hub"
)

print(cache_dir)
\end{dhkcode}

다만 이 코드는 기본 경로를 계산할 뿐, 현재 Hugging Face 라이브러리가 실제로 사용하는 캐시 위치를 조회하는 코드는 아니다.

실제 캐시 위치는 관련 환경변수와 라이브러리 설정을 확인해야 한다.

\subsubsection{Hugging Face Hub 캐시의 내부 구조}

Hub 캐시는 일반적으로 모델 저장소와 버전을 구분할 수 있는 구조를 사용한다.

예를 들어 다음과 같은 형태다.

\begin{dhkdir}[in=2]
hub/
> models--distilbert--distilbert-base-uncased/
>> blobs/
>> refs/
>> snapshots/
\end{dhkdir}

각 디렉터리의 역할을 살펴보자.

\begin{dhktable}
|디렉터리|	역할|
|:---|:---|
|blobs/|	다운로드한 파일 내용 저장|
|refs/|	브랜치나 태그 등의 참조 정보|
|snapshots/	|특정 리비전의 파일 구성을 나타내는 경로|
\end{dhktable}

일반적인 캐시 구조에서는 \zsh{snapshots}의 파일들이 \zsh{blobs}에 저장된 파일을 가리키는 심볼릭 링크로 구성될 수 있다.

환경에 따라 링크를 사용할 수 없는 경우에는 저장 방식이 달라질 수 있다.

여기서 Git과 비슷한 구조가 등장하는 이유는 모델 저장소도 버전이 존재하기 때문이다.

예를 들어 같은 모델 이름이라도 저장소가 업데이트되면 서로 다른 리비전의 파일이 존재할 수 있다.

따라서 모델 이름만으로 모든 실행 환경이 완전히 고정되는 것은 아니다.

\subsubsection{\zsh{HF_HOME}과 \zsh{HF_HUB_CACHE}}

이제 환경변수를 연결해 보자.

Hugging Face의 캐시 위치는 환경변수를 통해 설정할 수 있다.

대표적인 환경변수는 다음과 같다.

\begin{dhktable}
|환경변수|	역할|
|:---|:---|
|HF_HOME|	Hugging Face 관련 로컬 데이터의 기본 위치|
|HF_HUB_CACHE|	Hugging Face Hub 저장소 캐시 위치|
|HF_TOKEN|	Hub 인증에 사용할 토큰|
|HF_HUB_OFFLINE|	Hub의 네트워크 접근을 비활성화하는 설정|
\end{dhktable}

\begin{dhknote}
\begin{dhksmry}[centered]
HF_HOME
\end{dhksmry}
\end{dhknote}

예를 들어 macOS 터미널에서 다음과 같이 설정할 수 있다.

\begin{dhkzshcode}
export HF_HOME="$HOME/huggingface-data"
\end{dhkzshcode}

그러면 Hugging Face 라이브러리는 해당 위치를 기본 디렉터리로 사용할 수 있다.

Hub 캐시의 기본 위치도 일반적으로 그 아래에 구성된다.

\begin{dhkdir}[in=6]
huggingface-data/
> hub/
\end{dhkdir}

\begin{dhknote}
\begin{dhksmry}[centered]
HF_HUB_CACHE
\end{dhksmry}
\end{dhknote}

Hub 캐시 위치만 별도로 지정할 수도 있다.

\begin{dhkzshcode}
export HF_HUB_CACHE="$HOME/model-cache"
\end{dhkzshcode}

이 경우 Hub 모델 저장소 캐시의 위치를 별도로 관리할 수 있다.

두 변수를 동시에 지정했다면 \zsh{HF_HUB_CACHE}가 Hub 캐시 위치를 직접 결정하는 설정이다.

\begin{dhknote}
\begin{dhksmry}[centered]
Python에서 설정하기
\end{dhksmry}
\end{dhknote}

이전에 배운 \zsh{os.environ}을 사용할 수도 있다.

\begin{dhkcode}
import os

os.environ["HF_HOME"] = "/path/to/huggingface-data"
\end{dhkcode}

하지만 여기서 중요한 점이 있다.

Hugging Face 관련 라이브러리를 import하기 전에 환경변수를 설정하는 것이 안전하다.

일부 라이브러리는 import 과정에서 환경변수를 읽어 내부 설정을 초기화하기 때문이다.

예를 들어 다음과 같은 순서가 좋다.

\begin{dhkcode}
import os

os.environ["HF_HOME"] = "/path/to/huggingface-data"

from transformers import AutoModel
\end{dhkcode}

반대로 이미 라이브러리를 import한 뒤 환경변수를 변경하면,

변경한 설정이 현재 프로세스의 라이브러리 동작에 즉시 반영되지 않을 수 있다.

이것이 7장에서 배웠던 환경변수와 Python 객체의 초기화 시점이 연결되는 부분이다.

\subsubsection{환경변수와 \zsh{.env} 파일}

이번에는 \zsh{.env} 파일을 살펴보자.

먼저 중요한 사실부터 짚고 넘어간다.

\zsh{.env} 파일은 Python이나 운영체제가 자동으로 읽어 환경변수로 등록해 주는 특별한 파일이 아니다.

일반적으로 애플리케이션이나 별도의 도구가 \zsh{.env} 파일을 읽고 환경변수를 설정한다.

예를 들어 프로젝트에 다음 파일이 있다고 가정하자.

\begin{dhkdir}[in=6]
transformer-project/
> .env
\end{dhkdir}

내용

\begin{dhkzsh}
HF_HOME=/Users/user/huggingface-data
HF_TOKEN=your_token_here
\end{dhkzsh}

이 파일은 단순한 텍스트 파일이다.

Python이 실행되었다고 해서 자동으로 다음 값이 설정되는 것은 아니다.

\begin{dhkcode}
os.environ["HF_TOKEN"]
\end{dhkcode}

\zsh{.env}를 사용하려면 별도의 로딩 과정이 필요하다.

\begin{dhknote}
\begin{dhksmry}[centered]
python-dotenv 사용하기
\end{dhksmry}
\end{dhknote}

대표적인 방법은 \zsh{python-dotenv} 라이브러리를 사용하는 것이다.

\begin{dhkzsh}
python -m pip install python-dotenv
\end{dhkzsh}

그다음에 아래처럼 작성할 수 있다.

\begin{dhkcode}
from pathlib import Path
from dotenv import load_dotenv

PROJECT_ROOT = Path(__file__).resolve().parent.parent

load_dotenv(
    PROJECT_ROOT / ".env"
)
\end{dhkcode}

이제 \zsh{.env} 파일의 설정을 현재 Python 프로세스의 환경변수로 읽어 들일 수 있다.

\begin{dhkcode}
import os

token = os.environ.get("HF_TOKEN")
\end{dhkcode}

여기서 중요한 구분이 있다.

\begin{dhkcode}
load_dotenv(...)
\end{dhkcode}

는 \zsh{.env} 파일의 값을 환경변수로 로딩하는 작업이고,

\begin{dhkcode}
os.environ.get("HF_TOKEN")
\end{dhkcode}

은 현재 프로세스의 환경변수에서 값을 조회하는 작업이다.

또한 \zsh{load_dotenv()}는 기본적으로 이미 존재하는 환경변수를 덮어쓰지 않는다.

필요한 경우 \zsh{override=True}를 지정할 수 있다.

\begin{dhknote}
\begin{dhksmry}[centered]
인증 토큰을 다룰 때의 주의사항
\end{dhksmry}
\end{dhknote}

Hugging Face 토큰은 인증 정보이므로 소스 코드에 직접 작성하거나 로그에 출력하지 않는 것이 좋다.

특히 다음 코드는 피해야 한다.

\begin{dhkcode}
print(os.environ["HF_TOKEN"])
\end{dhkcode}

그리고 \zsh{.env} 파일을 Git 저장소에 올리지 않도록 설정해야 한다.

\zsh{.gitignore}에 다음 항목을 추가할 수 있다.

\begin{dhkzsh}
.env
\end{dhkzsh}

다만 \zsh{.gitignore}는 아직 추적되지 않은 파일이 Git에 새로 추가되는 것을 방지하는 용도다.

이미 Git에서 추적 중인 파일을 자동으로 제거해 주지는 않는다.

\subsubsection{Hugging Face 인증과 공개 모델}

여기서 인증이 항상 필요한 것은 아니라는 점도 알아두자.

예를 들어 다음과 같은 공개 모델은 일반적으로 별도의 인증 토큰 없이 다운로드할 수 있다.

\begin{dhkcode}
from transformers import AutoTokenizer

tokenizer = AutoTokenizer.from_pretrained(
    "distilbert-base-uncased"
)
\end{dhkcode}

반면 접근 제한이 있는 모델이나 비공개 저장소에 접근하려면 인증이 필요할 수 있다.

토큰을 명시적으로 전달하는 방식도 있지만, 라이브러리 버전에 따라 인수 이름이 다를 수 있다.

예를 들어 최근 버전에서는 다음과 같은 형태를 사용할 수 있다.

\begin{dhkcode}
import os

from transformers import AutoModel

model = AutoModel.from_pretrained(
    "some-org/private-model",
    token=os.environ["HF_TOKEN"]
)
\end{dhkcode}

Hugging Face CLI로 로그인한 경우에는 저장된 인증 정보를 라이브러리가 사용할 수도 있으므로, 반드시 매번 토큰을 코드에 전달해야 하는 것은 아니다.

\subsubsection{오프라인에서 모델 사용하기}

이번에는 다운로드한 모델을 네트워크 없이 사용하는 방법을 알아보자.

앞에서 캐시가 존재한다고 했다.

그렇다면 이미 다운로드한 모델은 인터넷 연결 없이도 사용할 수 있을까?

필요한 파일이 로컬에 모두 존재한다면 가능하다.

대표적으로 다음 옵션을 사용할 수 있다.

\begin{dhkcode}
from transformers import AutoTokenizer, AutoModel

model_name = "distilbert-base-uncased"

tokenizer = AutoTokenizer.from_pretrained(
    model_name,
    local_files_only=True
)

model = AutoModel.from_pretrained(
    model_name,
    local_files_only=True
)
\end{dhkcode}

여기서

\begin{dhkcode}
local_files_only=True
\end{dhkcode}

는 로컬에 존재하는 파일만 사용하도록 요청하는 옵션이다.

필요한 파일이 로컬에 없으면 다운로드하는 대신 오류가 발생한다.

따라서 다음 두 상황을 구분해야 한다.

\begin{dhkdir}[in=6]
상황 A
> 필요한 파일이 로컬에 있음
> 오프라인 로딩 가능
\end{dhkdir}

\begin{dhkdir}[in=6]
상황 B
> 필요한 파일이 로컬에 없음
> 오프라인 로딩 실패
\end{dhkdir}

환경변수로 Hub의 오프라인 모드를 활성화할 수도 있다.

macOS/Linux:

\begin{dhkzsh}
export HF_HUB_OFFLINE=1
\end{dhkzsh}

다만 \zsh{HF_HUB_OFFLINE}은 Hub 라이브러리의 네트워크 접근을 제한하는 설정이므로, 애플리케이션 전체의 모든 네트워크 접근을 차단하는 것은 아니다.

\subsubsection{모델 캐시와 모델 저장 디렉터리의 차이}

여기서 혼동하기 쉬운 두 개념을 구분한다.

첫 번째는 Hugging Face가 관리하는 캐시다.

\begin{dhkzsh}
~/.cache/huggingface/hub/
\end{dhkzsh}

두 번째는 개발자가 직접 저장하는 모델 디렉터리다.

\begin{dhkdir}[in=6]
transformer-project/
>models/
>>my-model/
\end{dhkdir}

예를 들어 다음 코드를 보자.

\begin{dhkcode}
model.save_pretrained(
    "models/my-model"
)

tokenizer.save_pretrained(
    "models/my-model"
)
\end{dhkcode}

이 코드는 모델과 토크나이저를 지정한 디렉터리에 저장한다.

그다음 아래처럼 다시 불러올 수 있다.

\begin{dhkcode}
model = AutoModel.from_pretrained(
    "models/my-model"
)

tokenizer = AutoTokenizer.from_pretrained(
    "models/my-model"
)
\end{dhkcode}

두 방식의 차이를 정리하면 다음과 같다.

\begin{dhktable}
|구분	|Hugging Face 캐시|	save_pretrained() 디렉터리|
|:---|:---|:---|
|주요 목적|	다운로드 파일 재사용|	모델 결과물 저장 및 재사용|
|관리 주체|	Hugging Face 라이브러리|	개발자|
|위치	|캐시 설정에 따라 결정|	개발자가 지정|
|버전 구조|	Hub 리비전 기반|	개발자가 관리|
|프로젝트 산출물|	일반적으로 아님|	일반적으로 해당|
\end{dhktable}

특히 모델을 파인튜닝했다면 다운로드 캐시와 파인튜닝 결과물을 구분해서 관리하는 것이 좋다.

\subsubsection{environment.yml}과 실행 환경 재현

이제 Conda 환경으로 넘어가자.

우리는 2장에서 Conda 환경을 만들고 활성화하는 방법을 배웠다.

그런데 다른 컴퓨터에서 동일한 프로젝트를 실행하려면 어떻게 해야 할까?

프로젝트 소스 코드만 복사한다고 충분하지는 않다.

Python 버전과 필요한 라이브러리도 설치해야 한다.

이때 사용할 수 있는 것이 \zsh{environment.yml}이다.

예를 들어 다음과 같은 파일을 작성할 수 있다.

YAML :

\begin{dhkzsh}
name: transformer-project

channels:
  - conda-forge

dependencies:
  - python=3.10
  - pip
  - pandas
  - numpy
  - pip:
      - torch
      - transformers
      - datasets
      - accelerate
      - python-dotenv
\end{dhkzsh}

이 파일은 Conda 환경을 구성하기 위한 선언적인 설정 파일이다.

다음 명령으로 환경을 생성할 수 있다.

\begin{dhkzsh}
conda env create -f environment.yml
\end{dhkzsh}

그리고 활성화한다.

\begin{dhkzsh}
conda activate transformer-project
\end{dhkzsh}

여기서 중요한 점은 \zsh{environment.yml} 자체가 Python 인터프리터나 라이브러리를 포함하는 것은 아니라는 것이다.

환경을 구성하는 데 필요한 정보를 기술한 파일이다.

실제 패키지는 Conda와 pip가 설치한다.

\begin{dhknote}
\begin{dhksmry}[centered]
패키지 버전 고정
\end{dhksmry}
\end{dhknote}

위 예제에서는 라이브러리 버전을 명시하지 않았다.

그러면 설치 시점에 따라 서로 다른 버전이 설치될 수 있다.

따라서 재현성이 중요한 프로젝트에서는 다음과 같이 버전을 명시하는 것이 좋다.

\begin{dhkzsh}
dependencies:
  - python=3.10
  - pip
  - pip:
      - transformers==4.30.2
      - datasets==2.14.7
      - accelerate==0.20.3
\end{dhkzsh}

이 버전들은 이전에 공부했던 Transformer 학습 환경과 관련된 예시이다.

다만 실제 환경을 재구성할 때는 PyTorch, NumPy, Hugging Face Hub 등

다른 의존성의 버전 호환성도 함께 검증해야 한다.

특히 macOS Apple Silicon과 Windows에서는 사용할 수 있는 패키지 빌드가 다를 수 있다.

따라서 동일한 \zsh{environment.yml}을 사용하더라도 모든 운영체제에서 완전히 동일한 바이너리 환경이 만들어진다고 보장할 수는 없다.

\subsubsection{\zsh{sys.executable}로 실제 환경 검증하기}

여기서 이전에 배운 \zsh{sys.executable}이 다시 등장한다.

Conda 환경을 활성화했다고 생각해 보자.

\begin{dhkzsh}
conda activate transformer
\end{dhkzsh}

그다음 Python을 실행하면 일반적으로 활성화된 환경의 인터프리터가 선택된다.

하지만 VS Code, Jupyter, 터미널 설정 등에 따라 실제 사용 중인 인터프리터가 다를 수 있다.

따라서 다음 코드로 확인하는 것이 좋다.

\begin{dhkcode}
import sys

print(sys.executable)
\end{dhkcode}

예를 들어 macOS에서는 다음과 같은 경로가 나올 수 있다.

\begin{dhkzsh}
/Users/user/anaconda3/envs/transformer/bin/python
\end{dhkzsh}

Windows에서는 다음과 같을 수 있다.

\begin{dhkzsh}
C:\Users\user\anaconda3\envs\transformer\python.exe
\end{dhkzsh}

여기서 중요한 것은 환경 이름이 아니라 실제로 실행 중인 인터프리터의 경로다.

또한 설치된 라이브러리의 버전도 확인할 수 있다.

\begin{dhkcode}
import torch
import transformers

print(torch.__version__)
print(transformers.__version__)
\end{dhkcode}

이렇게 하면 현재 프로세스가 실제로 어떤 환경에서 실행되고 있는지 확인할 수 있다.

\subsubsection{Jupyter 커널과 Conda 환경의 관계}

Windows에서 Jupyter 커널 문제를 겪어본 경험이 있을 것이다.

Jupyter Notebook을 실행할 때는 다음 두 개념을 구분해야 한다.

\begin{dhksmry}
    1. Jupyter 서버를 실행하는 환경

    2. Notebook 코드가 실행되는 커널 환경
\end{dhksmry}

이 둘은 반드시 동일할 필요가 없다.

예를 들어 다음과 같은 구조가 가능하다.

\begin{dhkbranch}[in=8]
node main :: Jupyter Server
arrow main :: Notebook 실행 요청
node main :: Jupyter Kernel
handoff main -> j :: Python Interpreter
handoff j -> c :: Conda Environment
stop c
\end{dhkbranch}

따라서 터미널에서

\begin{dhkzsh}
conda activate transformer
\end{dhkzsh}

를 실행했다고 해서 이미 열려 있는 모든 Notebook이 자동으로 해당 환경을 사용하는 것은 아니다.

Notebook 내부에서 다음 코드를 실행해 확인해야 한다.

\begin{dhkcode}
import sys

print(sys.executable)
\end{dhkcode}

이 경로가 의도한 Conda 환겨으이 Python을 가리키는지 확인하는 것이다.

또한 커널에 필요한 라이브러리가 설치되어 있어야 한다.

예를 들어 \zsh{ipykernel}이 필요하다면 해당 환경에 설치할 수 있다.

\begin{dhkzsh}
conda activate transformer

python -m pip install ipykernel
\end{dhkzsh}

필요한 경우 커널을 등록할 수도 있다.

\begin{dhkzsh}
python -m ipykernel install \
    --user \
    --name transformer \
    --display-name "Python (transformer)"
\end{dhkzsh}

이렇게 등록하면 Jupyter에서 해당 커널을 선택할 수 있다.

여기서 \zsh{--name}은 커널의 내부 식별자이고,

\zsh{--display-name}은 사용자에게 표시되는 이름이다.

\subsubsection{실험 자동화와 \zsh{subprocess}}

이제 8장에서 배웠던 \zsh{subprocess}를 연결해 보자.

우리가 다음과 같은 학습 명령을 실행한다고 가정하자.

\begin{dhkzsh}
python -m src.train \
    --data ./data/train.csv \
    --epochs 10 \
    --batch-size 32 \
    --output ./checkpoints/model.pt
\end{dhkzsh}

그런데 배치 크기를 변경하면서 여러 실험을 수행하고 싶다면 어떻게 해야 할까?

예를 들어 다음 세 가지 실험을 수행한다고 생각해 보자.

\begin{dhktable}
|실험|	Batch Size	|Epoch|
|:---|:---|:---|
|Experiment 1|	16|	10|
|Experiment 2|	32|	10|
|Experiment 3|	64|	10|
\end{dhktable}

이때 매번 터미널에서 명령을 입력할 수도 있지만 Python으로 자동화할 수도 있다.

프로젝트에 다음 파일을 추가하자.

\begin{dhkdir}[in=6]
scripts/
>run_experimnets.py
\end{dhkdir}

코드는 다음과 같다.

\begin{dhkcode}
# scripts/run_experiments.py

import subprocess
import sys

from pathlib import Path


PROJECT_ROOT = Path(__file__).resolve().parent.parent

BATCH_SIZES = [16, 32, 64]


def main():

    for batch_size in BATCH_SIZES:

        command = [
            sys.executable,
            "-m",
            "src.train",
            "--data",
            "data/train.csv",
            "--epochs",
            "10",
            "--batch-size",
            str(batch_size),
            "--output",
            f"checkpoints/model_bs{batch_size}.pt"
        ]

        result = subprocess.run(
            command,
            cwd=PROJECT_ROOT,
            check=True
        )

        print(
            f"Batch size {batch_size}: "
            f"Exit code {result.returncode}"
        )


if __name__ == "__main__":
    main()
\end{dhkcode}

이 코드를 실행하면 세 개의 학습 실험이 순서대로 실행된다.

\begin{dhkzsh}
python scripts/run_experiments.py
\end{dhkzsh}

\begin{dhknote}
\begin{dhksmry}[centered]
sys.executable을 사용하는 이유
\end{dhksmry}
\end{dhknote}

다음 코드를 보자.

\begin{dhkcode}
sys.executable
\end{dhkcode}

이것은 현재 실행 중인 Python 인터프리터의 경로다.

따라서 자식 프로세스에서도 동일한 Python 인터프리터를 사용하도록 할 수 있다.

반면 다음처럼 작성하면

\begin{dhkcode}
"python"
\end{dhkcode}

Shell이나 운영체제의 실행 파일 검색 경로에 따라 다른 인터프리터가 선택될 가능성이 있다.

그래서 실험 자동화에서는 \zsh{sys.executable}을 사용하는 것이 유용하다.

\begin{dhknote}
\begin{dhksmry}[centered]
cwd=PROJECT_ROOT의 역할
\end{dhksmry}
\end{dhknote}

다음 코드를 보자.

\begin{dhkcode}
cwd=PROJECT_ROOT
\end{dhkcode}

이것은 자식 프로세스의 CWD를 프로젝트 루트로 지정한다.

따라서 자식 프로세서으ㅔ서 다음 상대 경로를 사용하면

\begin{dhkzsh}
data/train.csv
\end{dhkzsh}

프로젝트 루트 기준으로 해석된다.

또한 \zsh{python -m src.train}이 프로젝트 루트에서 실행되므로

\zsh{src} 패키지를 찾는 데도 도움이 된다.

\begin{dhknote}
\begin{dhksmry}[centered]
check=True의 역할
\end{dhksmry}
\end{dhknote}

다음 코드는

\begin{dhkcode}
check=True
\end{dhkcode}

자식 프로세스가 0이 아닌 종료 코드를 반환했을 때 \zsh{CalledProcessError}를 발생시킨다.

따라서 학습 실험이 실패하면 오류를 무시한 채 다음 실험으로 진행하는 것을 방지할 수 있다.

물론 실무에서는 실패한 실험을 기록하고 나머지 실험을 계속 실행하도록 설계할 수도 있다.

그것은 실험 관리 정책에 따라 결정하면 된다.

\subsubsection{실험 로그 관리하기}

이제 학습 로그를 파일로 저장해 보자.

다음처럼 프로젝트에 로그 디렉터리를 생성할 수 있다.

\begin{dhkcode}
log_dir = PROJECT_ROOT / "logs"

log_dir.mkdir(
    parents=True,
    exist_ok=True
)
\end{dhkcode}

그리고 각 실험의 로그 파일 경로를 만들자.

\begin{dhkcode}
log_path = log_dir / f"train_bs{batch_size}.log"
\end{dhkcode}

이제 다음처럼 실행할 수 있다.

\begin{dhkcode}
with log_path.open("w", encoding="utf-8") as log_file:

    result = subprocess.run(
        command,
        cwd=PROJECT_ROOT,
        stdout=log_file,
        stderr=subprocess.STDOUT,
        check=True
    )
\end{dhkcode}

여기서 \zsh{stdout=log_file}은 자식 프로세스의 표준 출력을 파일에 연결한다.

그리고

\begin{dhkcode}
stderr=subprocess.STDOUT
\end{dhkcode}

는 표준 오류 출력도 같은 대상으로 연결한다.

따라서 정상 출력과 오류 메시지를 하나의 로그 파일에 저장할 수 있다.

결과적으로 다음과 같은 구조가 만들어진다.

\begin{dhkdir}[in=6]
transformer-project/
>checkpoints/
>>model_bs16.pt
>>model_bs32.pt
>>model_bs64.pt
>logs/
>>train_bs16.log
>>train_bs32.log
>>train_bs64.log
\end{dhkdir}

여기서 중요한 점은 모델 파일과 로그 파일이 서로 다른 목적을 가진다는 것이다.

모델 파일은 학습 결과의 파라미터를 저장하고, 로그 파일은 실행 과정에서 발생한 메시지를 기록한다.

\subsubsection{실험 재현성(Reproducibility)}

이제 마지막으로 실험 재현성을 살펴보자.

머신러닝에서 재현성은 매우 중요한 개념이다.

그런데 단순히 난수 시드만 고정한다고 모든 실험이 재현되는 것은 아니다.

예를 들어 다음과 같이 작성할 수 있다.

\begin{dhkcode}
import random
import numpy as np
import torch

seed = 42

random.seed(seed)
np.random.seed(seed)
torch.manual_seed(seed)
\end{dhkcode}

이 코드는 Python, NumPy, PyTorch의 주요 난수 생성기를 초기화한다.

하지만 실제 학습 결과에는 다음과 같은 요소들도 영향을 줄 수 있다.

\begin{dhksmry}
    - 데이터셋 버전

    - 모델 초기 가중치

    - 라이브러리 버전

    - 하드웨어 종류

    - 연산의 결정성

    - 데이터 샘플링 순서

    - 학습 설정

    - 체크포인트의 상태
\end{dhksmry}

따라서 실험 재현성을 확보하려면 다음 정보들을 함께 관리하는 것이 좋다.

\begin{dhktable}
|구분|	기록할 정보|
|:---|:---|
|코드	|Git commit|
|환경	|Python 및 라이브러리 버전|
|데이터|	데이터셋 버전 또는 식별자|
|모델|	모델 이름과 리비전|
|설정|	Learning rate, batch size, epochs|
|난수|	Random seed|
|장치|	CPU, CUDA, MPS|
|결과|	Checkpoint, metrics, logs|
\end{dhktable}

특히 Hugging Face 모델을 사용할 때는 모델 이름뿐 아니라 리비전도 기록하는 것이 좋다.

예를 들어 다음처럼 작성할 수 있다.

\begin{dhkcode}
model = AutoModel.from_pretrained(
    model_name,
    revision=revision
)
\end{dhkcode}

여기서 \zsh{revision}에 특정 커밋 해시를 지정하면 해당 리비전의 모델 파일을 요청할 수 있다.

모델 저장소가 이후에 변경되더라도 동일한 리비전을 참조할 수 있다는 장점이 있다.

다만 정확한 재현성을 위해서는 토크나이저, 데이터셋, 라이브러리 환경도 함께 고정해야 한다.

\subsubsection{전체 프로젝트 실행 구조}

이제 1~9장에서 배운 내용을 하나의 구조로 정리해 본다.

\begin{dhkdiagram}[in=4]
node user (1,2) :: User
node shell (2,2) :: Shell
node conda (3,2) :: Conda Environment
node pp (4,2) :: Python Process
node argv (5,1) :: sys.argv
node env (5,2) :: os.environ
node cwd (5,3) :: Path.cwd()
node argp (6,1) :: argparse
node ev (6,2) :: Environment//Variables
node fp (6,3) :: File Paths
node train (7,2) :: train.py
node dset (8,1) :: Dataset
node conf (8,2) :: Config
node model (8,3) :: Model
node tens (9,1) :: Tensor
node param (9,2) :: Parameters
node weight (9,3) :: Weights
node tl (10,2) :: Training Loop
node check (11,1) :: Checkpoint
node metr (11,2) :: Metrics
node logs (11,3) :: Logs
node fs (12,2) :: File System

edge user -> shell :: CLI Command
edge shell -> conda :: conda activate transformer
edge conda -> pp :: Python Interpreter
split pp -> argv, env, cwd
edge argv -> argp
edge env -> ev
edge cwd -> fp
merge argp, ev, fp -> train 
split train -> dset, conf, model 
edge dset -> tens
edge conf -> param
edge model -> weight
merge tens, param, weight -> tl
split tl -> check, metr, logs
merge check, metr, logs -> fs
\end{dhkdiagram}

그리고 실험 자동화까지 포함하면 다음 구조가 추가된다.

\begin{dhkbranch}[in=8]
node main :: run_experiments.py
arrow main :: subprocess.run()
node main :: Child Python Process
arrow main :: python -m src.train
node main :: Training
branchline main -> c :: Checkpoint
stop c
branchline main -> m :: Metrics
stop m
branchline main -> l :: Logs
stop l
arrow main
node main :: Exit Code
arrow main
node main :: Parent Python Process
\end{dhkbranch}

이것이 지금까지 배운 Python 실행 환경의 전체 구조다.

\subsubsection{전체 커리큘럼 최종 정리}

\begin{dhktable}
|장|학습 주제|핵심 개념|
|:---|:---|:---|
|1장|Python 실행 환경과 외부 시스템 개요|sys, os, pathlib, 실행 환경과 외부 시스템의 관계|
|2장|파일 시스템과 경로|Path, CWD, 절대·상대 경로, pathlib|
|3장|경로의 기준점과 커맨드라인 실행|Path.cwd(), __file__, sys.path, Interpreter, Shell, sys.argv, argparse|
|4장|Python 모듈·패키지와 Import 시스템|module, package, sys.path, python -m, sys.modules, namespace, binding|
|5장|파일 입출력|Text, Binary, CSV, JSON|
|6장|URL과 네트워크|HTTP, requests, 다운로드|
|7장|환경변수|os.environ, PATH, PYTHONPATH|
|8장|프로세스|subprocess, stdin/stdout/stderr|
|9장|ML/DL 프로젝트 통합|학습 실행, 체크포인트, 캐시, 자동화|
\end{dhktable}

처음 이 커리큘럼을 구성한 이유는 Python의 기능을 하나씩 외우는 데 있지 않았다.

\begin{dhksmry}
    Python 프로그램이 어떤 환경에서 시작되고,

    외부 데이터를 어떻게 가져오며,

    객체를 어떻게 구성하고,

    다른 프로세스와 어떻게 상호작용하고,

    최종 결과를 어떻게 저장하는지 하나의 실행흐름으로 이해하는 것
\end{dhksmry}

이 목표였다.

그 목표에 필요한 핵심 개념들을 이번 9장까지 모두 연결하였다.